\section{Conclusion}
\label{sec:conclusion}

This article has surveyed verbal reinforcement learning through a single organizing question: where a natural-language signal enters the decision problem and what it modifies at the moment it takes effect, yielding the three pillars of language as grounding signal, as deliberative feedback, and as learning signal. The four contributions promised in \S\ref{sec:intro} were delivered as follows. The two-criterion fence and the language-augmented sequential decision framework of \S\ref{sec:framework} give VRL a definition tight enough to exclude plain instruction following, vanilla reinforcement learning from human feedback, and language-conditioned reinforcement learning, while the taxonomy of \S\ref{sec:taxonomy} assigns every signal by its point of effect and resolves hybrid cases, such as reward-code generation and experiential memory, by explicit decision rules rather than by fiat. These rules are turned outward in the prescriptive practitioner guide of \S\ref{sec:practitioner}, which keys the choice among verbal-feedback mechanisms to task verifiability, feedback latency and cost, and the dominant risk each mechanism carries. The synthesis sections then followed verbal signals into the machinery that consumes them: the verbal reward stack (\S\ref{sec:rewardstack}), credit assignment at outcome, process, and token granularity (\S\ref{sec:credit}), and optimizers that act in language space rather than on weights (\S\ref{sec:optimizers}). Finally, the comparison with prior surveys and the research agenda, also in \S\ref{sec:practitioner}, position this account against its neighbors and state what remains to be done.

The payoff of the signal-centric view is consolidation. Questions ordinarily scattered across separate literatures, such as how much of a critique survives compression to a scalar, which feedback forms pair with which algorithm families, and when a verifier should displace a judge, become instances of one question about signal placement and fidelity. Looking ahead, we expect verbal feedback to become a first-class design primitive in agent architectures. Deployed agents will increasingly consume grounding, deliberative, and learning signals within a single trajectory, and the bottleneck will shift from generating feedback to verifying it; the infrastructure that shift demands, from feedback provenance to standardized quality evaluation to machine-readable tool interfaces, is the substance of the research agenda in \S\ref{sec:practitioner}.

\subsection{Limitations}
\label{subsec:limitations}

The survey's own limits should be stated plainly. Within each category we cover representative papers rather than an exhaustive enumeration, and the taxonomy is an organizing lens rather than a strict partition: several methods span roles, since verbal feedback can serve simultaneously as a grounding signal, a deliberative mechanism, and a learning signal, and our decision rules adjudicate placement without dissolving the underlying overlap. Coverage is restricted to methods in which verbal feedback is explicit and central to the agent's process; adjacent work where language plays an auxiliary role is treated only in passing, and the definitional fence fixes the feedback modality to text, so spoken and richer multimodal feedback channels fall outside our scope even where embodied observations do not (\S\ref{sec:embodied}). The article is also a synthesis: it contributes no new benchmark or controlled empirical comparison, and where the literature offers none we have said so rather than inferred one.

Language-based feedback itself carries failure modes that this survey can name but not yet bound. A specification compiled from language can be executable yet mean the wrong thing, and policies exploit whatever a rubric or verifier leaves unspecified (\S\ref{sec:rewardstack}) \citep{mahmoud2026reward}. The consuming model can misread the signal: generative verifiers are fooled by superficial ``master keys'' \citep{zhao2025one}, and preference models reward convincing agreement over correctness \citep{sharma2024towards}. Grounding between verbal descriptions and environment state remains weak where it matters most: vision-language judges recognize objects but fail to understand task sequence \citep{wybitul2024vista}, and single-frame success detectors miss failures that live in the motion rather than the final state \citep{hwang2024motif}. On real-world transfer the surveyed record is positive but thin: Text2Reward, DrEureka, and RACER report successful real-robot deployments \citep{xie2024text2reward,ma2024dreureka,dai2024racer}, yet each is a demonstration on a specific platform, and we found no systematic study of when language-guided policies fail to transfer, so the breadth of transfer is an open question rather than an established property. Societal impact follows the same channels. The failures above already reach deployed settings, with sycophantic answer flips documented on medical-advice tasks \citep{fanous2025syceval} and injection attacks demonstrated against tool-calling, coding, and computer-use agents \citep{dziemian2026how}; in the high-stakes domains the field is entering (\S\ref{sec:intro}), a corrupted grounding signal harms both the system's users and the people affected by its outputs, who rarely chose the feedback channel and cannot audit it. Sections~\ref{sec:safety} and~\ref{sec:credit} discuss mitigation directions and the current state of formal guarantees; narrowing the gap between what the definitional fence admits and what theory can certify is, in our judgment, the field's most consequential open problem.
